La baixa resolució d'una partitura digitalitzada pot deteriorar pentagrames, símbols, dígits i
text. La superresolució permet generar una imatge més gran, però els models entrenats amb imatges
naturals poden produir detalls visualment plausibles que no estaven presents en el document. Aquest
Treball de Fi de Grau estudia quan aquestes tècniques aporten valor i quines precaucions exigeix el
seu ús professional en biblioteques, arxius, conservació, consulta o preparació editorial.

Es va auditar Sheet Music Benchmark (SMB) i es va reservar una mostra intacta per a la comparació
principal. Aquesta va utilitzar 64 pàgines de 64 obres independents i sis degradacions
sintètiques: factors d'ampliació dos i quatre, combinats amb severitat neta, moderada i forta. Perquè
la severitat fora comparable entre pàgines, el desenfocament es va normalitzar per la separació
entre línies de pentagrama. Es van comparar la interpolació bicúbica, la xarxa residual EDSR i el
transformador lleuger SwinIR mitjançant PSNR i SSIM sobre luminància, intervals aparellats per obra,
temps d'inferència i una revisió qualitativa preseleccionada d'errors de notació. La implementació
reproduïble es va desenvolupar en Python amb PyTorch, OpenCV i \texttt{uv}, i l'execució accelerada
es va realitzar en Kaggle sobre GPU Tesla T4.

Després de preservar la comparació principal, es va executar un estudi secundari d'ajust fi
d'EDSR. Es van definir rols interns per obra amb 45 grups d'entrenament, 13 de validació i 20 de
prova, sense solapament ni reutilització de les 64 obres principals. El model adaptat va superar
l'EDSR oficial en PSNR-Y i SSIM-Y per a les sis condicions, amb intervals aparellats que no van
creuar zero. Aquest resultat demostra adaptació dins d'SMB i degradacions coincidents, no
transferència a altres col·leccions ni correcció musical.

Com a extensió professional, el model adaptat es va aplicar sense reajustar-lo a dotze obres d'un
arxiu musical extern sota les mateixes degradacions. Els intervals van afavorir el model adaptat en
les sis comparacions PSNR-Y i en cinc de sis SSIM-Y. Huit derivats es van acceptar per a consulta,
tres amb reserves i un es va rebutjar; també es va implementar un demostrador local reversible.
Aquesta prova dona suport a la transferència entre col·leccions sota LR sintètica, no a la
restauració d'entrades LR naturals ni a una taxa universal de seguretat.

EDSR i SwinIR van superar la interpolació bicúbica en les dues mètriques primàries per a les sis
condicions. No obstant això, SwinIR només va millorar clarament EDSR en les condicions x2 neta, x2
moderada i x4 neta, i va ser entre 5,2 i 8,5 vegades més lent. La revisió visual va detectar símbols
alterats, corrupció de text i traços corbs o amb textura fins i tot quan augmentava la fidelitat
mitjana. Es conclou que EDSR ofereix el compromís més defensable entre fidelitat i cost per a generar
còpies de consulta assistides, mentre que SwinIR s'ha de reservar per a les condicions on el guany
compense el cost. Cap mètode queda validat com a restauració semàntica, substitut arxivístic o
entrada autònoma d'OMR: cal conservar l'original i verificar humanament cada derivat rellevant.
